\section{The second incompleteness theorem}
\label{sec:second-incompleteness-03}

\begin{theorem}[G\"odel's second incompleteness theorem]
\label{thm:godel-second}
Let \(T\) be a consistent computably enumerable extension of \(\mathsf Q\) whose provability
formula, chosen as in \cref{sec:second-incompleteness-01}, satisfies the
derivability conditions.  Then
\[
  T\notderives\mathsf{Con}(T).
\]
\textit{Status: complete proof.}
\end{theorem}

\begin{proof}
Suppose \(T\derives\mathsf{Con}(T)\).  From
\(\neg\mathsf{Prov}_T(\overline{\ulcorner\bottom_A\urcorner})\), assume
\(\mathsf{Prov}_T(\overline{\ulcorner\bottom_A\urcorner})\); negation
elimination yields \(\bottom\), and bottom elimination yields \(\bottom_A\).
Implication introduction therefore gives
\[
  T\derives
  \mathsf{Prov}_T(\overline{\ulcorner\bottom_A\urcorner})\to\bottom_A.
\]
The hypothesis of L\"ob's theorem now holds with \(\varphi=\bottom_A\).
Therefore \(T\derives\bottom_A\), contradicting the assumed consistency of
\(T\).
\end{proof}

\begin{warning}[Scope of the second theorem]
The conclusion is not ``\(T\) cannot prove any sentence expressing any notion
of consistency.''  It concerns the particular provability formula chosen in
\cref{sec:second-incompleteness-01} and the corresponding sentence
\(\mathsf{Con}(T)\).  Other deliberately weak predicates can fail the
derivability conditions.
\end{warning}

\begin{example}[Relative consistency is not excluded]
A stronger theory \(T'\) may prove \(\mathsf{Con}(T)\); the second theorem
places no prohibition on such a relative result.  For example, sufficiently
strong standard set theories verify the existence of the standard natural
number structure and its satisfaction of the axioms of \(\PA\), from which
they formalize \(\mathsf{Con}(\PA)\).  This is a claim about what the stronger
theory proves, not a claim that its consistency alone implies
\(\mathsf{Con}(T)\).  The theorem says that a consistent \(T\) cannot prove its
own fixed consistency sentence under the stated hypotheses.
\end{example}
